\documentclass{article}
\usepackage[utf8]{inputenc}
\usepackage[spanish]{babel}
\usepackage[right=2cm,left=3cm,top=2cm,bottom=2cm,headsep=0cm,footskip=0.5cm]{geometry}


\title{Pauta Presentación}
\author{Marina Fabregat Expósito}
\date{\today}

\begin{document}
\maketitle

\section{Introducción:}
\subsection{Justificación o motivación del trabajo}

Motivaciones:

\begin{itemize}
    \item El curso en la upf.
    \item Es un tema que me interesa. 
    \item Gente de mi alrededor me ha motivado, familia, conocidos, etc.
\end{itemize}

Objetivos:

\begin{itemize}
    \item Me gustaría trabajar en seguridad en la red y he visto este trabajo como una introducción a ello. 
    \item Conocer más sobre la criptografía clásica.
    \item Saber como funciona \LaTeX$\,$ ya que es una herramienta muy útil en el ámbito académico. 
\end{itemize}
\subsection{Hipótesis}
Una de las dos hipótesis es si soy capaz de comprender los diferentes métodos de criptografía clásica. Hay operaciones matemáticas algo más difíciles de las que he aprendido en lo que llevo de bachillerato. (Por eso la hipótesis)

Otra de las hipótesis es ver si soy capaz de saber utilizar \LaTeX, que es un sistema de composición de textos, como word, pero orientado más en el ámbito científico. 

\section{Desarrollo del trabajo}
\subsection{Metodología y trabajo de campo}
Con ayuda de documentos y información que me ha proporcionado mi tutor he recogido toda la información necesaria para desarrollar los puntos del trabajo teórico. He ido añadiendo aclaraciones en algunos lugares del trabajo para que me sea más fácil comprender el temario. 

El marco práctico de este trabajo se basa en una serie de ejercicios donde he cifrado, descifrado y he hecho el criptoanálisis a una serie de textos y frases con diferentes métodos que he explicado en el marco teórico.

El \LaTeX $\,$ también forma parte del trabajo de campo del TR. Explicar en que consiste.

Después he organizado la información en diferentes bloques y la he pasado al ordenador con el recopilador de textos \LaTeX $\,$ que me ha sido muy útil porque las formulas matemáticas son más sencillas de escribir y el resultado es mucho mejor que el de un word. 

\subsection{Problemas y soluciones}

Criptografía: He tenido problemas con las matemáticas ya que en algunos casos he tenido que aprender unas nuevas operaciones y conceptos que no se explican en bachillerato. (Calcular el inverso modular de un número, por ejemplo). He tenido ayuda de diferentes webs que me daban el resultado y en algunos casos la explicación para obtenerlo. 

\LaTeX: Al principio era complicado adaptarse a este nuevo sistema, pero con ayuda de documentos, paginas web y mi tutor he podido adaptarme. Uno de los problemas que más me surgían al principio era que yo tenía una idea de lo que quería hacer pero no sabía como escribirla para que saliera tal y como yo quería. Por eso he empleado mucho tiempo en buscar en diferentes paginas web diversos métodos para conseguir mi objetivo. 
(Explicar algún que otro problema con \LaTeX, como el de poner código sin que el programa lo lea...)

\section{Conclusiones}

Las dos hipótesis que había formulado al principio de mi trabajo se han cumplido porque, a pesar de los inconvenientes y problemas, he conseguido comprender el funcionamiento de \LaTeX y los diferentes sistemas de cifra que he investigado, y esto se puede ver reflejado en los ejercicios del marco práctico y por la parte del \LaTeX en el propio trabajo, es decir como esta organizado y el formato de este.

\end{document}